\chapter{Diseño metodológico}
\label{cap:metodologia}

\section{Tipo y enfoque metodológico}
El Grupo~9.2 clasifica el estudio como \textbf{investigación tecnológica} en el marco de
\textit{design science}: se diseña un artefacto (modelo de alineación S/E/C $\rightarrow$
\AlignmentScore{}) que transforma entradas documentales en resultados verificables. No es un estudio
experimental sobre personas: la ``señal'' son archivos OpenAPI y extractos BIAN, no encuestas ni
sensores físicos (Clase~04; §1.9).

El plan (Guía~1) documenta el \textbf{diseño reproducible} del procedimiento. La construcción del
software y el experimento controlado quedan para Tesis~2 (§1.6), pero este capítulo anticipa
cómo se adquirirán los datos, qué unidad se medirá y qué factores se contrastarán.


\section{Procedimiento general}
El procedimiento sigue las cinco fases del modelo (Anexo~C), alineadas con OE1--OE4. La
\textbf{adquisición de artefactos} (Cap.~4 §4.3) precede a la normalización: contratos OpenAPI y
material BIAN se identifican, versionan y almacenan antes del \textit{parsing}. Este capítulo
documenta el diseño del procedimiento; su ejecución a escala queda fuera del alcance (§1.6).

\begin{enumerate}
  \item \textbf{Fase 1 --- Diseñar el modelo:} definir entradas, scores S/E/C, fórmula de
        agregación y umbrales de clasificación (Cap.~1 §1.4; Anexo~C).
  \item \textbf{Fase 2 --- Adquirir y normalizar artefactos (OE1):} obtener contratos OpenAPI y
        extractos BIAN (§4.3); parsear y producir representaciones comparables (tabular, grafo o RDF
        según factor elegido).
  \item \textbf{Fase 3 --- Emparejar estructuras (OE2):} calcular correspondencias y
        \StructuralScore{} (E).
  \item \textbf{Fase 4 --- Calcular similitud semántica (OE3):} obtener \SemanticScore{} (S).
  \item \textbf{Fase 5 --- Integrar y clasificar (OE4):} calcular \CoverageScore{} (C),
        \AlignmentScore{} y clasificación por tipologías (Tabla~\ref{tab:tipologias}).
\end{enumerate}

\subsection{Relación objetivos, fases y procedimiento}
\label{sec:objetivos-metodologia}

Esta subsección explicita cómo cada objetivo específico se desarrolla en el procedimiento. La
\Cref{tab:oe-metodologia} vincula OE, fase, acción metodológica y producto verificable.

\begin{table}[H]
  \centering
  \caption{Objetivos específicos vinculados al procedimiento metodológico}
  \label{tab:oe-metodologia}
  \small
  \begin{tabular}{@{}p{0.06\textwidth}p{0.14\textwidth}p{0.22\textwidth}p{0.24\textwidth}p{0.24\textwidth}@{}}
    \toprule
    \textbf{OE} & \textbf{Fase} & \textbf{Acción (verbo)} & \textbf{Entrada} & \textbf{Producto} \\
    \midrule
    --- &
    Fase 1 &
    \textbf{Diseñar} entradas, scores y umbrales &
    Requisitos del §1.4; Anexo C &
    Modelo S/E/C documentado \\
    \textbf{OE1} &
    Fase 2 &
    \textbf{Adquirir} y \textbf{normalizar} contrato y SD &
    OpenAPI + extracto BIAN (§4.3) &
    Representaciones comparables \\
    \textbf{OE2} &
    Fase 3 &
    \textbf{Emparejar} estructuras y calcular E &
    Representaciones OE1 &
    Matriz de correspondencias + E \\
    \textbf{OE3} &
    Fase 4 &
    \textbf{Calcular} similitud semántica S &
    Representaciones OE1 &
    Matriz semántica + S \\
    \textbf{OE4} &
    Fase 5 &
    \textbf{Integrar} S, E, C y clasificar &
    E, S, representaciones OE1 &
    C, \AlignmentScore{}, tipologías \\
    \bottomrule
  \end{tabular}
\end{table}

Cada fila corresponde a un hito del Capítulo~\ref{cap:objetivos} y a la matriz de consistencia
(Anexo~C). Las técnicas concretas por OE se detallan en la \Cref{sec:tecnicas-herramientas}.


La \Cref{fig:proc-estructural} ilustra la correspondencia jerárquica para E; la cadena completa
aparece en la \Cref{fig:cadena-objetivos} (Cap.~2).

\setcounter{figure}{6}
\begin{figure}[H]
  \centering
  \figCapSeven
  \caption{Correspondencia jerárquica para el cálculo de \StructuralScore{} (E). Fuente: elaboración
  propia; ejemplo ilustrativo.}
  \label{fig:proc-estructural}
\end{figure}

\subsection{Protocolo de validación empírica (fuera de alcance)}

Para una validación empírica posterior, el protocolo escrito contemplará: (1)~instrumento de captura
(repositorio de contratos); (2)~infraestructura (\textit{parsing}, almacenamiento);
(3)~condiciones de registro; (4)~tamaño y representatividad de la muestra; (5)~etiquetado de pares
OpenAPI$\leftrightarrow$BIAN (juicio experto o \textit{ground truth} documentado);
(6)~aspectos éticos/legales; (7)~evaluación del desempeño del sistema (precisión vs.\ referencia,
análisis de falsos positivos/negativos en matching).


\section{Fuentes de datos}
\subsection{Adquisición de datos}

Los datos no provienen de campo clínico ni de encuestas: son \textbf{artefactos digitales}
(contratos OpenAPI en YAML/JSON y material de referencia BIAN) que el procedimiento trata como
representaciones documentales comparables. Cada instancia debe identificarse de forma unívoca (p. ej.
\texttt{\{fuente\}\_\{service-domain\}\_\{version\}}).

\begin{table}[H]
  \centering
  \caption{Adquisición de datos --- diseño del experimento (Clase~04 / PI-11)}
  \label{tab:adquisicion-datos}
  \small
  \begin{tabular}{@{}p{0.26\textwidth}p{0.66\textwidth}@{}}
    \toprule
    \textbf{Elemento} & \textbf{En el Grupo~9.2} \\
    \midrule
    Sistema / entorno &
    Sector bancario; arquitecturas API + modelos de referencia BIAN \\
    Objeto de estudio &
    Contrato OpenAPI bancario completo (Cap.~1 §1.4) \\
    Unidad de análisis &
    Par comparable endpoint/\allowbreak schema $\leftrightarrow$ \Behavior{}/\BusinessObject{} \\
    Señal / dato &
    Archivos OpenAPI (YAML/JSON) + extracto del \ServiceDomain{} BIAN \\
    Instrumento de medición &
    Pipeline del modelo (normalización + cálculo S, E, C); \textit{parsing} OpenAPI \\
    Parámetros controlados &
    Versión OpenAPI~3.x; dominio SD; idioma de descripciones; herramienta de \textit{parsing} \\
    Parámetros no controlados &
    Calidad redaccional del contrato; madurez BIAN del emisor; completitud del SD publicado \\
    Factores (Cap.~2) &
    Esquema intermedio; estrategia de matching; técnica semántica; pesos $\alpha,\beta,\gamma$ \\
    Ratios (resultados) &
    \StructuralScore{}, \SemanticScore{}, \CoverageScore{}, \AlignmentScore{}, clasificación \\
    Formato almacenamiento &
    YAML/JSON versionado en repositorio; metadatos CSV/JSON por instancia \\
    Tamaño de muestra &
    Por definir en Tesis~2; ilustración actual: Anexo~F (Payment Initiation) \\
    \bottomrule
  \end{tabular}
\end{table}


\subsection{Criterios de inclusión (muestra empírica)}

Para una validación empírica posterior se considerarán contratos OpenAPI~3.x del sector bancario
emparejados con un \ServiceDomain{} BIAN declarado en la documentación del caso, con:

\begin{itemize}
  \item especificación parseable sin errores de sintaxis;
  \item al menos una operación y un esquema en \texttt{components/schemas};
  \item referencia BIAN disponible en formato legible (catálogo, extracto o documentación oficial);
  \item ausencia de datos personales identificables (Ley N°~29733): preferir ejemplos públicos,
        sintéticos o anonimizados (Anexo~F).
\end{itemize}

La \textbf{población teórica} abarca todas las instancias de medición posibles sobre pares
OpenAPI $\leftrightarrow$ BIAN en ese universo; la \textbf{muestra empírica} será el subconjunto de
contratos/pares seleccionados según representatividad de dominios (pagos, cuentas, préstamos, etc.),
sin fijar un piloto único obligatorio.


\section{Técnicas y herramientas}
Las técnicas se seleccionaron a partir del estado del arte (Cap.~3) y se vinculan a los factores
del Cap.~2:

\begin{itemize}
  \item \textbf{OE1:} \textit{parsing} OpenAPI~3.x; extracción de \textit{paths}, operaciones y
        esquemas; alineación con elementos BIAN en representación intermedia (Mainas et al., 2023
        como referencia de enriquecimiento semántico).
  \item \textbf{OE2:} \textit{schema matching} según taxonomía Shvaiko y Euzenat (2005): nombre,
        ruta, profundidad de esquema, operación HTTP.
  \item \textbf{OE3:} similitud semántica por embeddings y/o ontología (El Bayadh et al., 2015;
        Chandrasekaran y Mago, 2022).
  \item \textbf{OE4:} agregación ponderada $\alpha S + \beta E + \gamma C$ y clasificación por
        umbrales fijos (§1.4).
\end{itemize}

Herramientas previstas: Python para prototipo, validadores OpenAPI, bibliotecas de similitud textual;
detalle de implementación fuera de alcance de Guía~1.


\section{Métricas de evaluación}
Los ratios principales se definen operativamente en el Cap.~1 §1.4 y en el Anexo~D. Resumen:

\begin{itemize}
  \item \textbf{\StructuralScore{} (E):} grado de correspondencia estructural OpenAPI $\leftrightarrow$
        BIAN normalizado a $[0,1]$.
  \item \textbf{\SemanticScore{} (S):} similitud de significado entre términos y conceptos.
  \item \textbf{\CoverageScore{} (C):} proporción de elementos BIAN cubiertos por el contrato.
  \item \textbf{\AlignmentScore{}:} combinación ponderada; clasificación Alta ($\geq 0{,}80$),
        Media, Baja, Nula según §1.4.
\end{itemize}

Los indicadores por fase se detallan en la matriz de consistencia (Anexo~C) y en la
\Cref{tab:anexo-d}.


\section{Casos de estudio}
Los casos de estudio para validación empírica se elegirán por \textbf{representatividad de dominio},
no como piloto único del plan. El Anexo~F (\textit{Payment Initiation}) muestra el formato esperado de
entrada/salida: pares \texttt{POST /payment-orders} $\leftrightarrow$ \Behavior{} \texttt{Initiate},
etc.

Este plan documenta que cada caso aportará: (a)~identificador de instancia; (b)~contrato
OpenAPI; (c)~extracto BIAN; (d)~tabla de pares UA evaluados; (e)~scores S, E, C y
\AlignmentScore{} esperados como salida del artefacto. La ejecución y el tamaño de muestra
permanecen fuera del alcance (§1.6).


\section{Criterios de validez}

La validez del plan se sustenta en: (1)~\textbf{trazabilidad} problema $\rightarrow$ objetivos
$\rightarrow$ fases $\rightarrow$ productos (Anexo~C); (2)~\textbf{reproducibilidad} del
procedimiento documentado; (3)~\textbf{coherencia constructiva} entre definiciones Cap.~1, marco
Cap.~3 y operacionalización Anexo~D; (4)~\textbf{delimitación explícita} de alcance (§1.6).

\section{Amenazas a la validez}

\begin{itemize}
  \item \textbf{Interna:} ambigüedad léxica en contratos; errores de \textit{parsing}; sesgo en
        etiquetado experto de pares OpenAPI$\leftrightarrow$BIAN.
  \item \textbf{Externa:} muestra de contratos no representativa de todo el sector; versiones BIAN
        distintas entre emisores.
  \item \textbf{De constructo:} \AlignmentScore{} resume S/E/C; puede ocultar desalineaciones
        localizadas relevantes para negocio.
  \item \textbf{Conclusión:} generalizar scores de la muestra al universo de contratos bancarios
        requiere criterios de muestreo explícitos (§4.3).
\end{itemize}
